\chapter{ניהול הפרויקט}\label{ch:mgmt}

\section{עץ פירוק עבודה}

ראו איור~\ref{fig:wbs}.

\begin{figure}[!htbp]
\begin{ltrdiagram}
% wbs: TikZ / pgfplots source, built inside ltrdiagram (LTR)
\begin{tikzpicture}[font=\sffamily\scriptsize,
  root/.style={draw=acc2,fill=soft2,very thick,rounded corners=3pt,minimum width=3cm,minimum height=0.8cm,align=center},
  l1/.style={draw=acc,fill=soft,rounded corners=3pt,text width=2.3cm,minimum height=0.8cm,align=center},
  l2/.style={draw=gray!60,fill=white,rounded corners=2pt,text width=2.3cm,align=center,inner sep=2pt},
  ed/.style={thick,ink!50}]
\node[root] (r) at (7.5,0) {\en{\textbf{VeriLens 1.0}}};
\foreach \i/\n/\t/\a/\b/\c in {
 0/1/{\he{1. ניהול}}/{\he{תוכנית ולו״ז}}/{\he{סיכונים}}/{\he{דיווח לחוזה}},
 1/2/{\he{2. נתונים}}/{\he{מאגר אימון}}/{\he{מאגר כיול}}/\en{VL-Wild},
 2/3/{\he{3. מנוע אימות}}/{\he{שכבות 1 עד 3}}/{\he{שקלול וכיול}}/\en{API},
 3/4/{\he{4. מסוף ושילוב}}/{\he{מסוף העורך}}/{\he{תוספים למערכות תוכן}}/{\he{לוח בקרה}},
 4/5/{\he{5. תקן אתי}}/{\he{טיוטה וסדנאות}}/{\he{סבב ציבורי}}/{\he{נוהל הסמכה}},
 5/6/{\he{6. פיילוט}}/{\he{5 מערכות}}/{\he{הדרכה}}/{\he{מדידת יעדים}}}{
  \node[l1] (n\n) at (\i*2.75+0.6,-1.5) {\t};
  \draw[ed] (r.south)--++(0,-0.35)-|(n\n.north);
  \node[l2] (a\n) at (\i*2.75+0.6,-2.7) {\a};
  \node[l2] (b\n) at (\i*2.75+0.6,-3.4) {\b};
  \node[l2] (c\n) at (\i*2.75+0.6,-4.1) {\c};
  \draw[ed,thin] (n\n.south)--(a\n.north);
}
\end{tikzpicture}

\end{ltrdiagram}
\caption{עץ פירוק עבודה: שש חבילות עבודה, ולכל אחת שלושה תוצרים}\label{fig:wbs}
\end{figure}

\section{לוח זמנים}
הפרויקט נמשך \num{18} חודשים, מינואר \num{2027} עד יוני \num{2028}. הנתיב הקריטי עובר דרך בניית המאגרים, פיתוח הגלאים, הכיול, והקבלה הטכנית: עיכוב בכל אחד מהם מעכב את כל הפרויקט. התקן האתי מתנהל במקביל וכולל מרווח ביטחון של חודשיים (איור~\ref{fig:gantt}).

\begin{figure}[!htbp]
\begin{ltrdiagram}
% gantt: TikZ / pgfplots source, built inside ltrdiagram (LTR)
\begin{ganttchart}[
  x unit=0.62cm, y unit chart=0.52cm, vgrid={*1{draw=black!8}},
  title label font=\sffamily\scriptsize, bar label font=\sffamily\scriptsize,
  milestone label font=\sffamily\scriptsize\bfseries,
  bar/.append style={fill=acc!45,draw=acc}, bar height=0.55,
  group/.append style={fill=ink}, group label font=\sffamily\scriptsize\bfseries,
  milestone/.append style={fill=acc2,draw=acc2},
  link/.style={-{Stealth[length=4pt]},acc2,thick}]{1}{18}
\gantttitle{2027}{12}\gantttitle{2028}{6}\\
\gantttitlelist{1,...,18}{1}\\
\ganttgroup{\he{1. ניהול}}{1}{18}\\
\ganttbar[name=d1]{\he{2. מאגרי נתונים}}{1}{5}\\
\ganttbar[name=d2]{\he{מנוע: שכבות 1 עד 3}}{3}{9}\\
\ganttbar[name=d3]{\he{מנוע: שקלול וכיול}}{8}{11}\\
\ganttbar{\he{4. מסוף ושילוב}}{5}{12}\\
\ganttbar{\he{5. תקן אתי}}{2}{10}\\
\ganttmilestone[name=m1]{\he{קבלה טכנית}}{12}\\
\ganttbar[name=p]{\he{6. פיילוט}}{13}{18}\\
\ganttmilestone[milestone left shift=-0.5,milestone right shift=-0.5]{\he{קבלה מבצעית}}{18}
\ganttlink{d1}{d3}
\ganttlink{d3}{m1}
\ganttlink{m1}{p}
\end{ganttchart}

\end{ltrdiagram}
\caption{גאנט: הנתיב הקריטי עובר דרך המאגרים, הכיול והקבלה הטכנית}\label{fig:gantt}
\end{figure}

\section{אבני דרך}

ראו טבלה~\ref{tab:miles}.

\begin{table}[!htbp]\centering
\caption{אבני הדרך, התוצר בכל אחת והקשר לדוחות הקורס}\label{tab:miles}
\begin{tabularx}{\linewidth}{@{}l l Y@{}}
\toprule
\textbf{אבן דרך} & \textbf{חודש} & \textbf{תוצר מחייב} \\
\midrule
\en{M0}: אישור דוח מכין & \num{0} & מסמך זה, חתום \\
\en{M1}: מאגרים מוכנים & \num{5} & מאגר אימון, כיול ו-\en{VL-Wild} (אצל צד שלישי) \\
\en{M2}: דוח התקדמות & \num{7} & דוח התקדמות לפי תבנית הקורס; גלאי ראשון ותוצאות ביניים \\
\en{M3}: טיוטת תקן לציבור & \num{8} & \en{TRUTH-IMG} גרסה \num{0.9} לסבב הערות \\
\en{M4}: קבלה טכנית & \num{12} & תוצאות \req{T-01} עד \req{T-08} \\
\en{M5}: פיילוט פעיל & \num{13} & \num{5} מערכות חדשות מחוברות \\
\en{M6}: קבלה מבצעית ודוח סיום & \num{18} & תוצאות \req{T-09} עד \req{T-11}, \en{TRUTH-IMG} \num{1.0}, דוח סיום \\
\bottomrule
\end{tabularx}
\end{table}

\section{צוות ואחריות}

ראו טבלה~\ref{tab:raci}.

\begin{table}[!htbp]\centering
\caption{מטריצת \en{RACI}: \textbf{R} מבצע, \textbf{A} אחראי, \textbf{C} מתייעץ, \textbf{I} מעודכן}\label{tab:raci}
\begin{tabularx}{\linewidth}{@{}Y c c c c c@{}}
\toprule
\textbf{חבילת עבודה} & \textbf{מנכ״ל/ית} & \textbf{\en{CTO}} & \textbf{ראש/ת \en{ML}} & \textbf{ראש/ת אתיקה} & \textbf{מוצר} \\
\midrule
\num{1}. ניהול & A/R & C & I & I & C \\
\num{2}. מאגרי נתונים & I & A & R & C & I \\
\num{3}. מנוע אימות & I & A & R & I & C \\
\num{4}. מסוף ושילוב & I & A & C & C & R \\
\num{5}. תקן אתי & A & I & I & R & C \\
\num{6}. פיילוט & A & C & C & R & R \\
\bottomrule
\end{tabularx}
\end{table}

תפקידי הטבלה הם תפקידים בחברה המתוכננת; בשלב הדוח המכין שני המייסדים, גיא ופז, ממלאים את כולם. הצוות בשיאו מונה \num{14} אנשים: \num{6} מהנדסי תוכנה, \num{3} חוקרי למידת מכונה, \num{2} בודקי אימות (מודיעין גלוי), מעצב/ת מוצר, ראש/ת אתיקה עם רקע עיתונאי, ומנכ״ל/ית.

\section{בקרת שינויים}\label{sec:change}
כיוון שהדוח הוא חוזה, כל שינוי בתוצר, בקריטריון קבלה, בגרף מוגדר או במועד אבן דרך מחייב בקשת שינוי כתובה. הבקשה כוללת: מה משתנה, למה, ההשפעה על לו״ז, תקציב וסיכונים, ואישור בכתב של המרצה. שינויים ללא אישור אינם מחייבים את המרצה בבדיקה.
